\documentclass[11pt,letterpaper]{article}
\usepackage[T1]{fontenc}
\usepackage[utf8]{inputenc}
\usepackage{lmodern}
\usepackage[margin=1in,headheight=14pt]{geometry}
\usepackage{amsmath,amssymb,amsthm,mathtools,bm}
\usepackage{microtype,booktabs,array,longtable,enumitem}
\usepackage{xcolor}
\usepackage{fancyhdr}
\usepackage{hyperref}
\usepackage{bookmark}
\hypersetup{colorlinks=true,linkcolor=blue!45!black,citecolor=blue!45!black,
 urlcolor=blue!45!black,pdftitle={Thue--Morse Mellin Renormalization and Gamma-Tower Uniqueness},
 pdfauthor={ChatGPT},pdfsubject={Research extension with proofs and reproducible computations}}
\pagestyle{fancy}
\fancyhf{}
\fancyhead[L]{\small Thue--Morse Mellin renormalization}
\fancyhead[R]{\small Research article}
\fancyfoot[C]{\thepage}
\renewcommand{\headrulewidth}{0.3pt}
\setlength{\parskip}{4pt}
\setlength{\parindent}{1.2em}
\setlength{\emergencystretch}{2em}
\allowdisplaybreaks[2]
\numberwithin{equation}{section}
\newtheorem{theorem}{Theorem}[section]
\newtheorem{proposition}[theorem]{Proposition}
\newtheorem{lemma}[theorem]{Lemma}
\newtheorem{corollary}[theorem]{Corollary}
\theoremstyle{definition}
\newtheorem{definition}[theorem]{Definition}
\newtheorem{example}[theorem]{Example}
\theoremstyle{remark}
\newtheorem{remark}[theorem]{Remark}
\newcommand{\N}{\mathbb N}
\newcommand{\R}{\mathbb R}
\newcommand{\C}{\mathbb C}
\newcommand{\Z}{\mathbb Z}
\newcommand{\E}{\mathbb E}
\newcommand{\W}{\mathcal W}
\newcommand{\CM}{\mathrm{CM}}
\newcommand{\dd}{\,\mathrm d}
\newcommand{\eps}{\varepsilon}
\DeclareMathOperator{\dist}{dist}
\DeclareMathOperator{\supp}{supp}
\DeclareMathOperator{\Rea}{Re}
\DeclareMathOperator{\diam}{diam}
\newcommand{\coef}[2]{[#1^{#2}]}
\newcommand{\Dprime}[1]{\partial_s D(#1,a)}
\title{\vspace{-1em}\textbf{Thue--Morse Mellin Renormalization}\\[0.35em]
\Large Fabius-controlled derivative asymptotics,\\
positive dyadic solutions, and gamma-tower uniqueness}
\author{A research extension to the ProveIt Thue--Morse atlas}
\date{4 September 2026}

\begin{document}
\maketitle
\thispagestyle{empty}
\begin{abstract}
The signed Thue--Morse sequence has a positive Laplace product
$Q(t)=\prod_{j\ge0}(1-e^{-2^jt})$. We pair it with the Fabius uniform-sum
Laplace transform
$H(z)=\prod_{j\ge1}(1-e^{-z/2^j})/(z/2^j)$.
Their product satisfies the exact equation $\W(2t)=\W(t)/t$.
This elementary cancellation gives a positive lognormal-type Mellin kernel
and an all-orders expansion for the derivatives of the entire shifted
Thue--Morse Dirichlet function at its negative-integer zeros.
The coefficients are finite expressions in Bernoulli numbers and the shift
parameter. A global Taylor estimate yields a uniform dyadic
$q$-Gevrey remainder and a Gaussian-in-order bound after optimal truncation.

We also classify every completely monotone solution of the basic dyadic
equation by a positive measure on one multiplicative period. This produces
many positive gamma towers with the same normalization. More strongly,
distinct towers can have identical formal high-order expansions to every
order; their discrepancy is beyond all those corrections. Nevertheless,
the single boundary condition $-a f'(a)\to1$ as $a\downarrow0$ selects the
canonical solution and its full differential tower.
Complete proofs, explicit coefficients, positive Hankel consequences, and
independently evaluated numerical checks are provided. Classical ingredients
are separated from the results derived here; historical priority of the
new formulations is not asserted as established.
\end{abstract}
\vspace{0.35em}
\noindent\textbf{Keywords.} Thue--Morse sequence; Fabius function; entire Dirichlet
series; Mellin transform; completely monotone function; moment problem;
log-periodicity; beyond-all-orders asymptotics.

\clearpage
{\linespread{0.94}\small\setlength{\parskip}{0pt}\tableofcontents}
\clearpage

\section{The questions and the main answers}\label{sec:intro}

Let $s_2(n)$ denote the number of ones in the binary expansion of $n$, and put
\begin{equation}\label{eq:epsilon}
 \eps_n=(-1)^{s_2(n)},\qquad n=0,1,2,\ldots.
\end{equation}
The first terms are $1,-1,-1,1,-1,1,1,-1,\ldots$.
Rather than attempting another catalogue of word-combinatorial properties,
we investigate two analytic questions that connect directly with the
Fabius--Rvachev program: how rapidly do the derivatives of the associated
entire Dirichlet function grow at negative integers, and what actually
characterizes its gamma-type hierarchy?

Throughout this article $a>0$ is a real shift, $\lambda=\log 2$, and
\begin{equation}\label{eq:Dintro}
 D(s,a)=\sum_{n\ge0}\frac{\eps_n}{(n+a)^s}
 \qquad(\Rea s>1)
\end{equation}
denotes the Dirichlet series and subsequently its entire continuation.
There are two different limiting parameters: $\rho\to\infty$ describes
\emph{high negative order}, whereas $a\downarrow0$ describes a
\emph{boundary of the shift variable}. They should not be conflated.

\subsection{High negative orders are governed by a positive dyadic kernel}
Define
\begin{equation}\label{eq:mainquantities}
 \begin{split}
 Q(t)&=\prod_{j\ge0}(1-e^{-2^jt}),\qquad
 H(z)=\prod_{j\ge1}\frac{1-e^{-z/2^j}}{z/2^j},\\
 \W(t)&=Q(t)H(t),\qquad
 C_\rho(a)=\int_0^\infty t^{-\rho-1}e^{-at}Q(t)\dd t.
 \end{split}
\end{equation}
Each removable value in the product for $H$ is understood by continuity.
For an integer $m\ge0$,
\begin{equation}\label{eq:derivativeintro}
 \partial_sD(-m,a)=(-1)^m m!\,C_m(a).
\end{equation}
The important cancellation is
\begin{equation}\label{eq:cocycleintro}
 \W(2t)=\frac{\W(t)}t.
\end{equation}
It yields a positive one-periodic amplitude $\Theta$ and coefficient
polynomials $c_k(a)$ such that
\begin{equation}\label{eq:expansionintro}
 \frac{C_\rho(a)}{2^{\rho(\rho-1)/2}\Theta(\rho)}
 \sim\sum_{k\ge0}c_k(a)2^{-k\rho}.
\end{equation}
The expansion is accompanied by an explicit-form remainder bound, not merely
formal coefficient matching. At integer orders $\Theta(m)=K_0$, where
\begin{equation}\label{eq:K0intro}
 K_0=\int_0^\infty \W(t)\frac{\dd t}{t}
   =0.75072198040787095920970791405063469\ldots.
\end{equation}
For example, at the symmetric shift $a=1/2$,
\begin{equation}\label{eq:centerintro}
 \frac{C_m(1/2)}{K_0\,2^{m(m-1)/2}}
 \sim 1-\frac{2^{-2m}}9
     +\frac{248}{2025}2^{-4m}
     -\frac{4806656}{2679075}2^{-6m}+\cdots.
\end{equation}
All odd corrections vanish. Section~\ref{sec:asymptotics} proves these
statements, including uniformity and optimal truncation.

\subsection{The dyadic equation has a large positive solution space}
A function $f:(0,\infty)\to\R$ is \emph{completely monotone} if it is smooth
and $(-1)^j f^{(j)}(a)\ge0$ for every $j\ge0$.
We completely describe the functions in this class satisfying
\begin{equation}\label{eq:basicintro}
 f(a)=f(a/2)-f((a+1)/2).
\end{equation}
They correspond exactly to positive, locally finite measures $\nu$ on
$(0,\infty)$ invariant under multiplication by two:
\begin{equation}\label{eq:classintro}
 f(a)=\int_0^\infty e^{-at}Q(t)\,\nu(\dd t),
 \qquad \nu(2B)=\nu(B).
\end{equation}
The canonical solution $C_0(a)$ corresponds to $\nu(\dd t)=\dd t/t$.
The normalization $f(1/2)=\log2$ fixes only the mass of one multiplicative
period, not the distribution of that mass. Even the addition of all formal
high-order coefficients need not restore uniqueness
(Section~\ref{sec:ambiguity}).

\subsection{A boundary condition does restore uniqueness}
Within the completely monotone class, the equation
\eqref{eq:basicintro} together with
\begin{equation}\label{eq:boundaryintro}
 \lim_{a\downarrow0}(-a f'(a))=1
\end{equation}
characterizes $f=C_0$ uniquely. The proof identifies the representing
measure from all dyadic boundary limits. The stronger, convenient condition
that $f(a)+\log a$ have a finite limit also suffices.
Together with the differential ladder, this selects the entire gamma tower.
These are precise characterization statements; they do not claim that an
arbitrary gamma-like function is uniquely determined without a specified
function class or boundary condition.

\subsection{Relation to the atlas and the literature}\label{subsec:provenance}
The supplied atlas~\cite{Atlas} develops finite products, repeated summation,
Fabius--Rvachev connections, an entire Mellin continuation, and an automatic
gamma hierarchy. Its ``Automatic Barnes hierarchy'' frontier asks for
normalizations yielding a uniqueness theorem. The present work gives a
complete answer within the completely monotone class, and also exhibits a
specific obstruction to using high-order asymptotics alone.

The analytic study of Thue--Morse Dirichlet series predates this article by
decades; relevant works include Allouche--Cohen~\cite{AC},
Allouche--Mend\`es France--Peyri\`ere~\cite{AMP}, and
Allouche's later survey~\cite{Allouche2015}. Explicit linear-combination
identities give a complementary direction~\cite{Toth}.
In particular, the simplicity of the negative-integer zeros is
\emph{not} offered as a new result.
The Fabius probability and product background is also classical; see
Arias de Reyna~\cite{Arias}.
Log-periodic effects in Mellin analysis and moment indeterminacy of
lognormal-type laws have established precedents~\cite{FGD,Heyde}.

The contributions developed and proved here are the quantitative
renormalization theorem with its global remainder, the measure
classification and boundary characterization, and the construction of
distinct normalized towers sharing every formal high-order correction.
The exact product pairing and its subsidiary consequences organize their
proofs. These are candidate research contributions, not a claim of verified
historical priority. The targeted comparison did not locate the same main
statements in the inspected atlas or accessible literature, but it was not
an exhaustive search. In particular, the author's link to
Alkauskas's 2001 preprint~\cite{Alkauskas} could not be retrieved; no claim
about its detailed contents is made. This manuscript was prepared with
ChatGPT; it has not been peer-reviewed or formalized in Lean.

\section{From binary signs to an entire Mellin transform}\label{sec:background}

\subsection{The product and elementary endpoint bounds}
The two binary identities
\begin{equation}\label{eq:binary}
 \eps_{2n}=\eps_n,\qquad \eps_{2n+1}=-\eps_n
\end{equation}
show by induction that
\begin{equation}\label{eq:finiteproduct}
 \sum_{n=0}^{2^m-1}\eps_n z^n
   =\prod_{j=0}^{m-1}(1-z^{2^j}).
\end{equation}
For $|z|<1$, both sides converge as $m\to\infty$, so
\begin{equation}\label{eq:Qgenerating}
 Q(t)=\sum_{n\ge0}\eps_n e^{-nt}\qquad(t>0).
\end{equation}
Although the sum has signs, the product is strictly positive.

\begin{lemma}[Bounds needed for all Mellin operations]\label{lem:Qbounds}
The product $Q$ is positive and smooth on $(0,\infty)$ and satisfies
\begin{equation}\label{eq:Qscale}
 Q(t)=(1-e^{-t})Q(2t).
\end{equation}
For every integer $N\ge0$,
\begin{equation}\label{eq:Qflat}
 0<Q(t)\le 2^{N(N-1)/2}t^N\qquad(t>0),
\end{equation}
and
\begin{equation}\label{eq:Qinfty}
 0\le1-Q(t)\le\frac{e^{-t}}{1-e^{-t}}.
\end{equation}
In particular, $Q$ is smaller than any fixed power at zero and approaches
one exponentially at infinity.
\end{lemma}
\begin{proof}
Uniform convergence of the product and its differentiated logarithms on
compact subintervals of $(0,\infty)$ follows from the exponential growth of
$2^j$. Its factors are positive, and the sum of their deficits is finite.
The scaling equation follows by shifting the index.
For \eqref{eq:Qflat}, bound the first $N$ factors by $2^jt$, and the rest
by one. For \eqref{eq:Qinfty}, use
$1-\prod_j(1-u_j)\le\sum_j u_j$ for $0\le u_j\le1$, then
$2^j\ge j+1$.
\end{proof}

\subsection{Continuation, real zeros, and positive derivative data}
\begin{proposition}[Entire continuation and real zeros]\label{prop:entire}
For $a>0$, the function
\begin{equation}\label{eq:MellinD}
 M(s,a)=\int_0^\infty t^{s-1}e^{-at}Q(t)\dd t,
 \qquad D(s,a)=\frac{M(s,a)}{\Gamma(s)},
\end{equation}
is entire in $s$ and agrees with \eqref{eq:Dintro} where $\Rea s>1$.
Its real zeros are exactly $0,-1,-2,\ldots$, and every one is simple.
Moreover, \eqref{eq:derivativeintro} holds.
\end{proposition}
\begin{proof}
Fix a compact set of $s$-values. At zero, choose $N$ in
\eqref{eq:Qflat} large enough to dominate
$t^{\Rea s-1}|\log t|^jQ(t)$ for every derivative order $j$ under
consideration. At infinity, $e^{-at}$ dominates any such power and
logarithm. Differentiation under the integral is therefore locally uniform,
which proves that $M$ is entire.

For $\Rea s>1$, termwise integration of \eqref{eq:Qgenerating} is justified
by the corresponding absolutely convergent unsigned sum. It gives the
usual gamma integral for each $(n+a)^{-s}$.
For real $s$, the integral $M(s,a)$ is strictly positive.
The reciprocal gamma function has simple zeros precisely at the
nonpositive integers, with
\[
 \left.\frac{\mathrm d}{\mathrm ds}\frac1{\Gamma(s)}
 \right|_{s=-m}=(-1)^m m!.
\]
Multiplication by $M(s,a)$ proves every assertion.
\end{proof}

It is useful to keep $\rho$ real rather than immediately restrict it to an
integer. The positive quantities $C_\rho$ in
\eqref{eq:mainquantities} exist for every real $\rho$, although the
asymptotic statements below will use $\rho\ge0$.
For every integer $\ell\ge0$,
\begin{equation}\label{eq:cmderivatives}
 (-1)^\ell\partial_a^\ell C_\rho(a)
  =\int_0^\infty t^{\ell-\rho-1}e^{-at}Q(t)\dd t>0.
\end{equation}
Thus every $C_\rho$ is completely monotone. In particular,
\begin{equation}\label{eq:ladder}
 C_m'(a)=-C_{m-1}(a)\quad(m\ge1).
\end{equation}
Using \eqref{eq:Qscale} and substituting $t=2u$ gives
\begin{equation}\label{eq:Crhodyadic}
 C_\rho(a)=2^\rho\left\{C_\rho(a/2)-C_\rho((a+1)/2)\right\}.
\end{equation}
All of these identities follow from absolutely convergent integrals;
regularized sums are unnecessary here.

\subsection{The Fabius uniform-sum transform}
Let $U_1,U_2,\ldots$ be independent uniform random variables on $[0,1]$, and
set
\begin{equation}\label{eq:X}
 X=\sum_{j\ge1}2^{-j}U_j.
\end{equation}
This converges almost surely, takes values in $[0,1]$, and has mean $1/2$.
The distribution is the standard dyadic uniform-sum realization underlying
the Fabius function. Explicitly fixing \eqref{eq:X} avoids a possible
factor-of-two ambiguity among Fabius and Rvachev conventions.
Independence gives
\begin{equation}\label{eq:Hprob}
 H(z)=\E e^{-zX}
     =\prod_{j\ge1}\psi(z/2^j),
 \qquad \psi(z)=\frac{1-e^{-z}}z,\quad\psi(0)=1.
\end{equation}
Boundedness of $X$ proves that $H$ is entire. Equivalently, the product
converges locally uniformly since $\psi(z/2^j)=1+O_K(2^{-j})$ on every
compact set $K$.

The distributional identity $X\overset d=(U+X')/2$, with $X'$ an independent
copy, gives a direct connection with the usual functional equation.
If $F(x)=\Pr(X\le x)$ is extended by zero to $x\le0$ and by one to
$x\ge1$, convolution with the uniform law gives
\[
 F'(x)=2\{F(2x)-F(2x-1)\}.
\]
Consequently $F'(x)=2F(2x)$ on $0<x<1/2$.
For every $N\ge2$, the sum of the first $N$ uniform terms has a
$C^{N-2}$ spline density; convolution with the remaining probability law
preserves that regularity. Since $N$ is arbitrary, the density is smooth.
Only the entire transform \eqref{eq:Hprob} is needed below.

\begin{lemma}[Scaling, symmetry, and singularity distance]\label{lem:H}
The function $H$ satisfies
\begin{equation}\label{eq:Hscale}
 H(2z)=\psi(z)H(z),\qquad H(-z)=e^zH(z).
\end{equation}
Its zeros are exactly the nonzero integer multiples of $4\pi i$, with
multiplicities prescribed by the individual factors. The nearest zeros to
zero, $\pm4\pi i$, are simple. In $|z|<4\pi$,
\begin{equation}\label{eq:logH}
 \log H(z)=-\frac z2+\sum_{j\ge1}b_jz^{2j},\qquad
 b_j=\frac{B_{2j}}{2j(2j)!(4^j-1)}.
\end{equation}
Here the logarithm has value zero at the origin.
\end{lemma}
\begin{proof}
The first identity is an index shift. The second follows either from
$X\overset d=1-X$ or directly from $\psi(-z)=e^z\psi(z)$.
Each factor has zeros at $z=2^{j+1}\pi i k$, $k\ne0$.
Away from those zeros, the convergence of the logarithmic tail shows that
the product is nonzero. This proves the description of all zeros and their
nearest distance.

The Bernoulli expansion is obtained by integrating the power series for
$\frac12\coth(z/2)-1/z$:
\[
 \log\psi(z)=-\frac z2+
 \sum_{j\ge1}\frac{B_{2j}}{2j(2j)!}z^{2j},\qquad |z|<2\pi.
\]
Substitute $z/2^\ell$ and sum the resulting geometric series over
$\ell\ge1$. The nearest factor singularity moves the resulting radius to
$4\pi$, yielding \eqref{eq:logH}.
\end{proof}

\section{An exact lognormal-type kernel}\label{sec:kernel}

\subsection{The product cancellation and its envelope}
\begin{theorem}[Exact Mellin renormalization kernel]\label{thm:kernel}
Let $\W(t)=Q(t)H(t)$. Then $\W>0$ and
\begin{equation}\label{eq:Wscale}
 \W(2t)=t^{-1}\W(t),\qquad
 \W(2^jv)=2^{-j(j-1)/2}v^{-j}\W(v)\quad(j\in\Z).
\end{equation}
The function
\begin{equation}\label{eq:A}
 A(y)=2^{(y^2-y)/2}\W(2^y)
\end{equation}
is positive, smooth, and one-periodic. Hence
\begin{equation}\label{eq:lognormalW}
 \W(e^x)=\exp\!\left(-\frac{x^2}{2\lambda}+\frac x2\right)A(x/\lambda).
\end{equation}
There are constants $0<A_-\le A_+<\infty$ such that the expression on the
right with $A$ replaced by $A_-$ or $A_+$ is respectively a lower or an
upper bound for $\W(e^x)$, for every real $x$.
\end{theorem}
\begin{proof}
From \eqref{eq:Qscale} and \eqref{eq:Hscale},
\[
 \W(2t)=\frac{Q(t)}{1-e^{-t}}\,
         \frac{1-e^{-t}}tH(t)=\frac{\W(t)}t.
\]
Iteration gives the second identity for positive integers $j$, and solving
backwards gives it for negative integers as well.
Substitution into \eqref{eq:A} gives $A(y+1)=A(y)$.
Positivity and smoothness follow from those of the two factors.
The positive minimum and finite maximum of $A$ on $[0,1]$ give the bounds.
\end{proof}

The word ``lognormal-type'' refers to the Gaussian in $x=\log t$.
It does \emph{not} assert that $\W(t)\dd t/t$ is itself exactly a lognormal
law: the periodic factor $A$ remains present.

\subsection{The Mellin recurrence and phase amplitude}
Define
\begin{equation}\label{eq:Kdef}
 K(s)=\int_0^\infty t^{s-1}\W(t)\dd t\qquad(s\in\C).
\end{equation}
The two-sided Gaussian envelope proves local uniform absolute convergence
for every $s$, including every differentiated integral. Thus $K$ is entire
and is strictly positive on the real axis.

\begin{proposition}[Exact moments and a periodic amplitude]\label{prop:K}
For $s\in\C$ and $k\in\Z$,
\begin{equation}\label{eq:Kshift}
 K(s+1)=2^{s+1}K(s),\qquad
 K(s+k)=2^{ks+k(k+1)/2}K(s).
\end{equation}
In particular, $K(k)=2^{k(k+1)/2}K_0$ for every integer $k$.
The function
\begin{equation}\label{eq:Thetadef}
 \Theta(\rho)=2^{-\rho(\rho-1)/2}K(-\rho)
\end{equation}
is entire, one-periodic, and positive for real $\rho$.
\end{proposition}
\begin{proof}
Substitute $t=2u$ in \eqref{eq:Kdef} and use \eqref{eq:Wscale}:
$K(s)=2^sK(s-1)$. This yields the first recurrence and then its integer
iterates. Substituting the recurrence in \eqref{eq:Thetadef} proves
periodicity. The other statements have already been justified by the
integral.
\end{proof}

\subsection{Why the Mellin modulation is extremely small}
The periodic function $A$ has an absolutely convergent Fourier series
\begin{equation}\label{eq:Afourier}
 A(y)=\sum_{\ell\in\Z}a_\ell e^{2\pi i\ell y},\qquad
 a_\ell=\int_0^1 A(y)e^{-2\pi i\ell y}\dd y.
\end{equation}
For example, two integrations by parts already give a summable bound on
its Fourier coefficients. Integrating each Fourier mode against the
Gaussian in \eqref{eq:lognormalW} gives
\begin{equation}\label{eq:Kfourier}
 K(s)=\sqrt{2\pi\lambda}\,
 e^{\lambda(s+1/2)^2/2}
 \sum_{\ell\in\Z}a_\ell
 e^{-2\pi^2\ell^2/\lambda}e^{2\pi i\ell(s+1/2)}.
\end{equation}
The identity holds initially for real $s$, where absolute convergence
justifies all operations, and then for complex $s$ by analytic continuation.
Consequently,
\begin{equation}\label{eq:Thetafourier}
 \Theta(\rho)=\sqrt{2\pi\lambda}\,2^{1/8}
 \sum_{\ell\in\Z}(-1)^\ell a_\ell
 e^{-2\pi^2\ell^2/\lambda}e^{-2\pi i\ell\rho}.
\end{equation}
Mellin transformation suppresses the $\ell$th oscillatory mode by a
Gaussian factor in $\ell$. Section~\ref{sec:numerics} finds relative
variations of $\Theta$ around $10^{-18}$. These numerical observations are
not used as a proof that $\Theta$ is nonconstant; none of our main theorems
requires that additional assertion.

\subsection{Two exact expansions from the same periodic factor}
The identity $Q=\W/H$, together with \eqref{eq:logH}, gives the convergent
small-$t$ correction formula
\begin{equation}\label{eq:logQexact}
 \begin{split}
 \log Q(t)={}&-\frac{(\log t)^2}{2\lambda}+\frac12\log t
       +\log A(\log_2t)\\
 &+\frac t2-\sum_{j\ge1}b_jt^{2j},\qquad 0<t<4\pi.
 \end{split}
\end{equation}
At the other end, expanding $-\log Q(t)$ in exponentials gives
\begin{equation}\label{eq:logHlarge}
 \begin{split}
 \log H(t)={}&-\frac{(\log t)^2}{2\lambda}+\frac12\log t
       +\log A(\log_2t)\\
 &+\sum_{n\ge1}\frac{2^{v_2(n)+1}-1}{n}\,e^{-nt},\qquad t>0,
 \end{split}
\end{equation}
where $v_2(n)$ is the exponent of two dividing $n$.
Indeed, in $\sum_{j\ge0}\sum_{r\ge1}e^{-r2^jt}/r$, the coefficient of
$e^{-nt}$ is $n^{-1}\sum_{j=0}^{v_2(n)}2^j$.
All terms are nonnegative, so the rearrangement is justified by Tonelli's
theorem. Thus the same periodic amplitude links a Bernoulli correction at
one endpoint to a convergent exponential correction at the other.

\section{All-orders derivative growth with a uniform remainder}\label{sec:asymptotics}

\subsection{Coefficient polynomials}
Define
\begin{equation}\label{eq:Rdef}
 R_a(z)=\frac{e^{-az}}{H(z)}=\sum_{k\ge0}r_k(a)z^k,
 \qquad c_k(a)=2^{k(k+1)/2}r_k(a).
\end{equation}
The Taylor series has radius exactly $4\pi$, since the numerator cannot
cancel the simple poles at $\pm4\pi i$.
Put
\begin{equation}\label{eq:ellcoeffs}
 L_1=\frac12-a,\qquad L_{2j}=-b_j\ (j\ge1),\qquad
 L_{2j+1}=0\ (j\ge1).
\end{equation}
Then $R_a(z)=\exp(\sum_{\ell\ge1}L_\ell z^\ell)$ near zero, and a finite,
nonrecursive coefficient formula is
\begin{equation}\label{eq:finitecoeff}
 c_k(a)=2^{k(k+1)/2}
 \sum_{\substack{j_1,\ldots,j_k\ge0\\
                  j_1+2j_2+\cdots+kj_k=k}}
 \prod_{\ell=1}^k\frac{L_\ell^{j_\ell}}{j_\ell!},
 \qquad c_0=1.
\end{equation}
For computation, the equivalent recurrence
\begin{equation}\label{eq:coeffrec}
 n r_n=\sum_{\ell=1}^n\ell L_\ell r_{n-\ell},\qquad r_0=1,
\end{equation}
is often more economical. It is a convenience, not a prerequisite to the
closed expression \eqref{eq:finitecoeff}.

With $b=a-1/2$, the first coefficients are
\begin{align}
 c_0&=1,& c_1&=-2b,\notag\\
 c_2&=4b^2-\frac19,&
 c_3&=-\frac{32}{3}b^3+\frac89b,\label{eq:firstcoeff}\\
 c_4&=\frac8{2025}(10800b^4-1800b^2+31),\notag\\
 c_5&=-\frac{256b}{2025}(2160b^4-600b^2+31),\notag\\
 c_6&=\frac{2048}{2679075}
       (3810240b^6-1587600b^4+164052b^2-2347).\notag
\end{align}
Differentiation and reflection give useful independent checks:
\begin{equation}\label{eq:coeffchecks}
 c_k'(a)=-2^k c_{k-1}(a),\qquad
 c_k(1-a)=(-1)^k c_k(a).
\end{equation}
For the second identity, use $R_{1-a}(z)=R_a(-z)$.

\subsection{Statement of the uniform theorem}
For a compact interval $I\subset(0,\infty)$ and $0<r<4\pi$, set
\begin{equation}\label{eq:tube}
 \mathcal T_r=\{z\in\C:\dist(z,[0,\infty))\le r\},\qquad
 M_{I,r}=\sup_{a\in I,\ z\in\mathcal T_r}|R_a(z)|.
\end{equation}
The finiteness of this constant is part of the proof below.
It is a mathematically specified bound; no numerical enclosure for
$M_{I,r}$ is claimed in the computational tables.

\begin{theorem}[Uniform all-orders expansion]\label{thm:asymptotics}
For every compact $I\subset(0,\infty)$, every $0<r<4\pi$, every integer
$N\ge1$, every $a\in I$, and every real $\rho\ge0$,
\begin{equation}\label{eq:maintheorem}
 \frac{C_\rho(a)}{2^{\rho(\rho-1)/2}\Theta(\rho)}
   =\sum_{k=0}^{N-1}c_k(a)2^{-k\rho}+\mathcal R_N(\rho,a),
\end{equation}
where
\begin{equation}\label{eq:mainremainder}
 |\mathcal R_N(\rho,a)|
 \le M_{I,r}\,r^{-N}2^{N(N+1)/2-N\rho}.
\end{equation}
In particular, for each fixed $N$ the error is uniformly
$O_{I,N}(2^{-N\rho})$ as $\rho\to\infty$.
\end{theorem}

The estimate is useful both with fixed $N$ and with $N$ increasing with
$\rho$. The latter use is the reason for keeping the dependence on $N$
explicit.

\subsection{Step one: exact finite-level renormalization}
\begin{lemma}\label{lem:finitelevel}
Write $\rho=m+\theta$, where $m=\lfloor\rho\rfloor$ and $0\le\theta<1$,
and set $h=2^{-m}$. Then
\begin{equation}\label{eq:exactrenorm}
 C_{m+\theta}(a)
 =2^{m(m-1)/2+m\theta}
   \int_0^\infty t^{-\theta-1}\W(t)R_a(ht)\dd t.
\end{equation}
\end{lemma}
\begin{proof}
Iteration of the product equation gives
\[
 Q(2^{-m}t)=Q(t)\prod_{j=1}^m(1-e^{-t/2^j}).
\]
Separate each factor into $(t/2^j)\psi(t/2^j)$ and use the product for $H$:
\begin{equation}\label{eq:Qfinitelevel}
 Q(2^{-m}t)=2^{-m(m+1)/2}t^mQ(t)\frac{H(t)}{H(2^{-m}t)}.
\end{equation}
Substituting the old variable $2^{-m}t$ in the definition of $C_{m+\theta}$
now gives \eqref{eq:exactrenorm}, with all powers of two displayed.
\end{proof}

The point of this formula is that the rapidly changing order $m$ has been
removed from the Mellin weight. It occurs only in the explicit prefactor
and in the small argument of $R_a$.

\subsection{Step two: a Taylor bound valid on the entire positive axis}
\begin{lemma}[Boundedness on a pole-free tube]\label{lem:tubebound}
For $I\subset(0,\infty)$ compact and $r<4\pi$, the constant $M_{I,r}$ in
\eqref{eq:tube} is finite. Consequently, for every real $x\ge0$,
\begin{equation}\label{eq:taylorglobal}
 \left|R_a(x)-\sum_{k=0}^{N-1}r_k(a)x^k\right|
  \le M_{I,r}r^{-N}x^N\qquad(a\in I).
\end{equation}
\end{lemma}
\begin{proof}
The tube contains no zero of $H$. Thus $R_a$ is jointly continuous and
uniformly bounded when $a\in I$ and $z$ stays in a compact portion of the
tube. It remains to control its unbounded right-hand end.

For real $x>0$, let
\[
 \mu_x=\frac{\E[Xe^{-xX}]}{\E[e^{-xX}]}=-\frac{H'(x)}{H(x)}.
\]
Logarithmic differentiation of the independent factors gives
\begin{equation}\label{eq:tiltedmean}
 \mu_x=\sum_{j\ge1}
 \left(\frac1x-\frac{2^{-j}}{e^{x2^{-j}}-1}\right).
\end{equation}
Each summand is the mean of a uniform variable on $[0,2^{-j}]$ under a
positive exponential tilt; it lies between zero and $2^{-j}/2$, and is
also at most $1/x$. The upper bound by the untilted mean follows because
the derivative of a tilted mean is the negative tilted variance. Splitting at $J=\lfloor\log_2x\rfloor$ for $x\ge1$
therefore gives
\begin{equation}\label{eq:tiltedbound}
 0\le\mu_x\le\frac{\log_2x+2}{x}.
\end{equation}
For $|y|\le r$,
\[
 |H(x+iy)-H(x)|
 \le\E[e^{-xX}|e^{-iyX}-1|]
 \le |y|\mu_xH(x).
\]
The rightmost multiplier is at most $1/2$ for sufficiently large $x$,
uniformly in $|y|\le r$. Thus $|H(x+iy)|\ge H(x)/2$ there.

Since $Q(x)\le1$, we have $H(x)=\W(x)/Q(x)\ge\W(x)$.
The lower bound in Theorem~\ref{thm:kernel}, with $a_0=\min I>0$, yields
\[
 |R_a(x+iy)|\le\frac{2}{A_-}
 \exp\!\left[-a_0x+\frac{(\log x)^2}{2\lambda}
                   -\frac12\log x\right]\longrightarrow0.
\]
The remaining part of the tube with negative real part is compact. This
proves $M_{I,r}<\infty$.

Every closed disc of radius $r$ centered on $[0,\infty)$ lies in the tube.
Cauchy's estimate gives
$|R_a^{(N)}(x)|/N!\le M_{I,r}r^{-N}$ for all $x\ge0$.
Taylor's integral remainder along the real interval from $0$ to $x$ gives
\eqref{eq:taylorglobal}. Notice that this argument does not require
$x<4\pi$; the radius restriction is on the complex discs used to control
derivatives, not on the real argument of the final estimate.
\end{proof}

\subsection{Step three: integration of the global remainder}
\begin{proof}[Proof of Theorem~\ref{thm:asymptotics}]
Insert the finite Taylor polynomial from \eqref{eq:taylorglobal} into
\eqref{eq:exactrenorm}. The $k$th term integrates to
$r_k(a)h^kK(k-\theta)$.
The shift identity gives
\begin{equation}\label{eq:momentratio}
 \frac{K(k-\theta)}{K(-\theta)}
   =2^{-k\theta+k(k+1)/2}.
\end{equation}
Also,
\begin{equation}\label{eq:leadingfactor}
 2^{m(m-1)/2+m\theta}K(-\theta)
   =2^{\rho(\rho-1)/2}\Theta(\rho).
\end{equation}
Dividing the finite sum by this leading factor yields
$c_k(a)2^{-k\rho}$ term by term.

Positivity of the weight $t^{-\theta-1}\W(t)$ permits direct integration
of the absolute remainder bound. The resulting relative error is at most
\[
 M_{I,r}r^{-N}h^N\frac{K(N-\theta)}{K(-\theta)}
 =M_{I,r}r^{-N}2^{N(N+1)/2-N\rho}.
\]
This proves the estimate without interchanging an infinite Taylor series
and an improper integral.
\end{proof}

\subsection{Divergence, optimal truncation, and centering}
The Taylor series for $R_a$ has radius $4\pi$, so
\begin{equation}\label{eq:rlimsup}
 \limsup_{k\to\infty}|r_k(a)|^{1/k}=\frac1{4\pi}.
\end{equation}
Along an infinite subsequence the additional factor
$2^{k(k+1)/2}$ forces $|c_k(a)|^{1/k}\to\infty$.
Hence the formal power series $\sum c_k(a)h^k$ has radius zero for every
$a>0$. Its use in \eqref{eq:maintheorem} is asymptotic, not convergent.
For any $r<4\pi$, Cauchy's estimate also gives
$|c_k(a)|\le M_{I,r}r^{-k}2^{k(k+1)/2}$.
We call this specific bound \emph{dyadic $q$-Gevrey growth}; the displayed
formula fixes the convention, rather than relying on differing definitions
of $q$-Gevrey order in the literature.

Let
\begin{equation}\label{eq:tau}
 \tau=\rho+\log_2r-\frac12.
\end{equation}
The base-two logarithm of the order-dependent factor in
\eqref{eq:mainremainder} is
$\frac12(N-\tau)^2-\frac12\tau^2$.
If $\tau\ge1$, choose a nearest integer $N\ge1$. Then
\begin{equation}\label{eq:optimal}
 |\mathcal R_N(\rho,a)|
 \le 2^{1/8}M_{I,r}\,2^{-\frac12(\rho+\log_2r-1/2)^2}.
\end{equation}
This is an achievable upper bound after a specified truncation, not a
claim that the true error has exactly this size or that the estimate is
asymptotically sharp.

When $a=1/2$, the function $R_{1/2}$ is even. Therefore all odd
coefficients vanish and the small scale in \eqref{eq:centerintro} is
$4^{-m}$ instead of $2^{-m}$. Fixed-order errors may accordingly be
improved by retaining the next vanishing odd coefficient before applying
the theorem. For example, the error after the displayed $2^{-6m}$ term is
$O(2^{-8m})$.

\subsection{Growth of the entire function and its gamma tower}
Define the gamma hierarchy by
\begin{equation}\label{eq:Gammadef}
 \Gamma_m^{\mathrm{TM}}(a)
   =\exp\bigl(\partial_sD(-m,a)\bigr)
   =\exp\bigl((-1)^m m!C_m(a)\bigr).
\end{equation}
Thus the expansion describes the \emph{signed logarithms} of these
functions. Exponentiating a truncated logarithmic expansion without
tracking the very large factor $m!2^{m(m-1)/2}$ would not automatically
produce a small relative error for $\Gamma_m^{\mathrm{TM}}$ itself.
This distinction matters at high order.

\begin{corollary}[Order and type]\label{cor:entiretype}
For every fixed $a>0$, the entire function $s\mapsto D(s,a)$ has order two
and type $\lambda/2$. Equivalently, if
$M_D(R)=\max_{|s|\le R}|D(s,a)|$, then
\begin{equation}\label{eq:entiretype}
 \limsup_{R\to\infty}\frac{\log M_D(R)}{R^2}=\frac{\log2}{2}.
\end{equation}
\end{corollary}
\begin{proof}
For $R\ge2$ and $|s|\le R$, split the Mellin integral at one:
\[
 |M(s,a)|\le C_R(a)+\Gamma(R)a^{-R}.
\]
Theorem~\ref{thm:asymptotics}, periodic boundedness of $\Theta$, and the
usual gamma integral estimate give a logarithmic upper bound
$\lambda R^2/2+O_a(R\log R)$.
The Weierstrass product for $1/\Gamma(s)$ gives
$\max_{|s|\le R}|1/\Gamma(s)|\le\exp(O(R\log R))$.
For completeness, split its product at $n=\lceil2R\rceil$:
the first part contributes $O(R\log R)$ and
$\log(1+s/n)-s/n=O(R^2/n^2)$ controls the tail.
This proves the required upper bound for $D$.

For the reverse bound, Cauchy's estimate on the circle of radius one about
$-m$ gives
\[
 M_D(m+1)\ge|\partial_sD(-m,a)|=m!C_m(a).
\]
The main theorem implies
$\log C_m(a)=\lambda m(m-1)/2+\log K_0+o(1)$.
Dividing by $(m+1)^2$ proves the matching lower limit superior.
A positive finite limit superior in \eqref{eq:entiretype} establishes the
stated order and type.
\end{proof}

\section{Positive moments, Hankel determinants, and invisible phases}\label{sec:moments}

\subsection{Strict positivity and quantitative Hankel growth}
The substitution $x=1/t$ in $C_m$ gives
\begin{equation}\label{eq:Cmoment}
 C_m(a)=\int_0^\infty x^m w_a(x)\dd x,
 \qquad w_a(x)=x^{-1}e^{-a/x}Q(1/x)>0.
\end{equation}
Every integer moment exists. Thus the following positivity concerns the
\emph{derivative data} of $D$, not an unrelated Hankel matrix formed
straight from the signs $\eps_n$.

\begin{proposition}[Strict Hankel positivity]\label{prop:Hankel}
For integers $n,\ell\ge0$,
\begin{equation}\label{eq:Hankelpositive}
 \det[C_{i+j+\ell}(a)]_{i,j=0}^n>0.
\end{equation}
In particular,
\begin{equation}\label{eq:Turan}
 C_m(a)C_{m+2}(a)>C_{m+1}(a)^2,
 \qquad
 \frac{C_m(a)C_{m+2}(a)}{C_{m+1}(a)^2}\longrightarrow2.
\end{equation}
More generally, for fixed $n$ as $m\to\infty$,
\begin{equation}\label{eq:Hankelasympt}
 \begin{split}
 \det[C_{m+i+j}(a)]_{i,j=0}^n\sim{}&K_0^{n+1}
 2^{\frac{n+1}{2}m(m-1)+mn(n+1)+\frac{n(n+1)(n-1)}3}\\
 &\times\prod_{0\le i<j\le n}(2^j-2^i).
 \end{split}
\end{equation}
The relative error in \eqref{eq:Hankelasympt} is $O(2^{-m})$ for fixed
$a,n$, and is $O(2^{-2m})$ at $a=1/2$.
\end{proposition}
\begin{proof}
The matrix in \eqref{eq:Hankelpositive} is the Gram matrix of
$1,x,\ldots,x^n$ for the strictly positive measure $x^\ell w_a(x)\dd x$.
A nonzero polynomial cannot vanish almost everywhere for this measure, so
the matrix is positive definite.

For the asymptotic determinant, apply Theorem~\ref{thm:asymptotics} to its
finitely many entries. Factor $K_0 2^{m(m-1)/2}$ from each row,
$2^{mi+i(i-1)/2}$ from row $i$, and
$2^{mj+j(j-1)/2}$ from column $j$.
The remaining matrix is $[2^{ij}+O(2^{-m})]$.
Its limiting determinant is the nonzero Vandermonde product
$\prod_{i<j}(2^j-2^i)$. Summing the row and column exponents gives
\eqref{eq:Hankelasympt}; the centered improvement follows from $c_1(1/2)=0$.
The limiting ratio in \eqref{eq:Turan} follows directly from the leading
terms of the three entries.
\end{proof}

\subsection{All integer Mellin moments can miss a periodic perturbation}
The kernel $\W$ gives a particularly transparent moment-indeterminate
family. This is a specific realization of a phenomenon familiar from the
classical lognormal moment problem~\cite{Heyde}, not a claim that moment
indeterminacy itself is new.

\begin{proposition}[Integer-moment invisibility]\label{prop:invisible}
Let $p:\R\to\R$ be bounded, measurable, and one-periodic. If
\begin{equation}\label{eq:pweightedzero}
 \int_0^\infty \W(t)p(\log_2t)\frac{\dd t}{t}=0,
\end{equation}
then, for every integer $k$,
\begin{equation}\label{eq:allmomentzero}
 \int_0^\infty t^{k-1}\W(t)p(\log_2t)\dd t=0.
\end{equation}
If in addition $p$ is nonzero on a set of positive measure and
$\|p\|_\infty\le1$, the probability measures
\begin{equation}\label{eq:stieltjesfamily}
 \frac{\W(t)}{K_0t}\{1+\eta p(\log_2t)\}\dd t,
 \qquad -1<\eta<1,
\end{equation}
are distinct and have the same moments $2^{k(k+1)/2}$ for every $k\in\Z$.
\end{proposition}
\begin{proof}
Multiplication of $\W(t)$ by $p(\log_2t)$ preserves the scaling equation
\eqref{eq:Wscale}. Its Mellin transform therefore satisfies the same shift
recurrence as $K$, by absolute integrability and the same substitution.
A zero at $s=0$ consequently gives zeros at all integers.
Positivity, normalization, and the moment formula for
\eqref{eq:stieltjesfamily} follow immediately. Nontriviality of $p$ makes
the measures distinct.
\end{proof}

Nonzero functions satisfying \eqref{eq:pweightedzero} are easy to produce:
subtract the weighted average from any nonconstant bounded periodic
function and rescale its sup norm. Later we will impose a second,
unweighted mean-zero condition to preserve a gamma normalization as well.

\section{Classification of completely monotone dyadic solutions}\label{sec:classification}

\subsection{The measure theorem}
We use the Bernstein representation theorem: a completely monotone
function on $(0,\infty)$ is the Laplace transform of a unique positive
Radon measure on $[0,\infty)$, with finite Laplace transform at every
positive argument. One reference stating both representation and uniqueness
is Aguech--Jedidi~\cite[Theorem~1]{AJ}.

\begin{theorem}[All positive dyadic solutions]\label{thm:classification}
A function $f:(0,\infty)\to\R$ is completely monotone and satisfies
\begin{equation}\label{eq:dyadicequation}
 f(a)=f(a/2)-f((a+1)/2)
\end{equation}
if and only if there is a positive Radon measure $\nu$ on $(0,\infty)$
with $\nu(2B)=\nu(B)$ such that
\begin{equation}\label{eq:frepresent}
 f(a)=\int_0^\infty e^{-at}Q(t)\,\nu(\dd t).
\end{equation}
The measure is unique. It is determined by an arbitrary finite positive
measure $\nu_0$ on the half-open fundamental interval $[1,2)$, and the
correspondence is explicitly
\begin{equation}\label{eq:shellrepresentation}
 f(a)=\int_{[1,2)}\sum_{j\in\Z}
          e^{-a2^jv}Q(2^jv)\,\nu_0(\dd v).
\end{equation}
The zero measure corresponds to the zero solution.
\end{theorem}
\begin{proof}
Suppose first that $f$ is completely monotone. Write
$f(a)=\int_{[0,\infty)}e^{-at}\,\mu(\dd t)$ by Bernstein's theorem.
For a measure $\mu$, denote by $\mu(2\,\dd u)$ the pushforward under
$t\mapsto t/2$, so it assigns mass $\mu(2B)$ to a Borel set $B$.
Equation \eqref{eq:dyadicequation} and uniqueness of Laplace transforms give
\begin{equation}\label{eq:mumeasure}
 \mu(\dd u)=(1-e^{-u})\mu(2\,\dd u).
\end{equation}
This also proves $\mu(\{0\})=0$.
Since $Q(u)>0$ for $u>0$, define
$\nu(\dd u)=Q(u)^{-1}\mu(\dd u)$. It is locally finite on $(0,\infty)$.
Using $Q(u)=(1-e^{-u})Q(2u)$ in \eqref{eq:mumeasure} yields
$\nu(\dd u)=\nu(2\,\dd u)$, exactly the stated invariance.
Uniqueness follows from uniqueness of $\mu$.

Conversely, let $\nu$ be positive and invariant. Its mass on $[1,2)$ is
finite by local finiteness. Partitioning $(0,\infty)$ into the disjoint
half-open intervals $[2^j,2^{j+1})$ gives \eqref{eq:shellrepresentation}.
At $j\to-\infty$, the bound \eqref{eq:Qflat} with arbitrarily large $N$
makes the shell sum converge, even after multiplication by any fixed
power of $t$. At $j\to+\infty$, the factor $e^{-at}$ dominates every
power. The convergence is locally uniform for $a>0$ and persists under
all $a$-derivatives. Thus \eqref{eq:frepresent} is finite and completely
monotone. The scaling equation of $Q$ and invariance of $\nu$ verify
\eqref{eq:dyadicequation} directly. Finally, a finite measure on $[1,2)$
extends uniquely by scaling to such a measure on $(0,\infty)$.
\end{proof}

This classification includes discrete, absolutely continuous, and singular
measures. It does not secretly assume that a density with respect to
$\dd t/t$ exists.

\subsection{Normalization fixes mass, not phase}
\begin{corollary}[The normalized solution simplex]\label{cor:normalization}
For the solution in Theorem~\ref{thm:classification},
\begin{equation}\label{eq:normalization}
 f(1/2)=\nu_0([1,2)).
\end{equation}
Consequently, the normalized solutions $f(1/2)=\lambda$ are exactly the
mixtures of
\begin{equation}\label{eq:phasefunctions}
 f_\theta(a)=\lambda\sum_{j\in\Z}
      e^{-a2^{j+\theta}}Q(2^{j+\theta}),\qquad 0\le\theta<1.
\end{equation}
Distinct phases give distinct functions. Their uniform phase average is
canonical:
\begin{equation}\label{eq:canonicalaverage}
 C_0(a)=\int_0^1 f_\theta(a)\dd\theta.
\end{equation}
\end{corollary}
\begin{proof}
The identity
\begin{equation}\label{eq:telescoper}
 e^{-t/2}Q(t)=Q(t)-Q(t/2)
\end{equation}
gives, for every fixed $v>0$,
\[
 \sum_{j\in\Z}e^{-2^{j-1}v}Q(2^jv)
 =\lim_{J\to\infty}\{Q(2^Jv)-Q(2^{-J-1}v)\}=1.
\]
Insert this into \eqref{eq:shellrepresentation} and use positivity to
interchange sum and integral. This proves \eqref{eq:normalization}.
The extremal measures with mass $\lambda$ are point masses
$\lambda\delta_{2^\theta}$; the representation is affine and injective,
so their functions are the extremal normalized solutions.
Distinct phases have disjoint multiplicative lattices, and uniqueness of
Laplace transforms distinguishes the functions.
For the last identity, substitute $t=2^{j+\theta}$ on each shell, so
$\dd t/t=\lambda\dd\theta$.
\end{proof}

In particular,
\begin{equation}\label{eq:canonicalnorms}
 C_0(1/2)=\lambda,\qquad C_0(1)=\lambda/2,
 \qquad\Gamma_0^{\mathrm{TM}}(1/2)=2.
\end{equation}
Every normalized solution also has $f(1)=\lambda/2$, because the equation
at $a=1$ reads $f(1)=f(1/2)-f(1)$.
Thus evaluating only at either of these two points conceals the
nonuniqueness. The computations below use $a=1/3$ instead.

\subsection{Nonuniqueness persists for entire differential towers}
Given an invariant positive measure $\nu$, define, for $\rho\in\R$,
\begin{equation}\label{eq:generalizedC}
 C_\rho^\nu(a)=\int_0^\infty t^{-\rho}e^{-at}Q(t)\,\nu(\dd t).
\end{equation}
The endpoint argument in the classification proof shows that these
integrals are finite and completely monotone for every fixed $\rho$.
For integer levels,
\begin{equation}\label{eq:nutowerlaws}
 (C_m^\nu)'=-C_{m-1}^\nu\quad(m\ge1),\qquad
 C_m^\nu(a)=2^m\{C_m^\nu(a/2)-C_m^\nu((a+1)/2)\}.
\end{equation}
Thus the phase freedom is compatible with every level of the differential
ladder and every dyadic equation. Moreover, $C_m^\nu(a)\to0$ as
$a\to\infty$, by dominated convergence with any fixed positive lower
bound on $a$. Vanishing at infinity does not eliminate this freedom.

\section{Universal coefficients and beyond-all-orders ambiguity}\label{sec:ambiguity}

\subsection{The same coefficient calculus works for every phase measure}
For a nonzero invariant positive measure $\nu$, set
\begin{equation}\label{eq:Knu}
 K_\nu(s)=\int_0^\infty t^s\W(t)\,\nu(\dd t),\qquad
 \Theta_\nu(\rho)=2^{-\rho(\rho-1)/2}K_\nu(-\rho).
\end{equation}
The Gaussian envelope on each shell proves that $K_\nu$ is entire and
positive on the real axis. Invariance and \eqref{eq:Wscale} imply
\begin{equation}\label{eq:Knurecurrence}
 K_\nu(s+1)=2^{s+1}K_\nu(s).
\end{equation}
Consequently $\Theta_\nu$ is one-periodic and positive on the real axis.

\begin{theorem}[Universal high-order coefficients]\label{thm:universal}
For every nonzero invariant positive measure $\nu$, the expansion
\begin{equation}\label{eq:universalexp}
 \frac{C_\rho^\nu(a)}{2^{\rho(\rho-1)/2}\Theta_\nu(\rho)}
 =\sum_{k=0}^{N-1}c_k(a)2^{-k\rho}+\mathcal R_N^\nu(\rho,a)
\end{equation}
has exactly the same coefficient polynomials as in
Theorem~\ref{thm:asymptotics}, and
\begin{equation}\label{eq:universalremainder}
 |\mathcal R_N^\nu(\rho,a)|
 \le M_{I,r}r^{-N}2^{N(N+1)/2-N\rho}.
\end{equation}
In particular, the relative remainder bound is independent of $\nu$.
\end{theorem}
\begin{proof}
Write $\rho=m+\theta$ as before. Scaling by $2^{-m}$ in the invariant
measure and using \eqref{eq:Qfinitelevel} gives
\[
 C_{m+\theta}^\nu(a)
 =2^{m(m-1)/2+m\theta}
   \int_0^\infty t^{-\theta}\W(t)R_a(2^{-m}t)\,\nu(\dd t).
\]
Apply the global Taylor estimate and replace the $K$-recurrence in the
proof of Theorem~\ref{thm:asymptotics} by
\eqref{eq:Knurecurrence}. Positivity permits the same absolute remainder
estimate. The moment ratios are identical, so no coefficient changes.
\end{proof}

\subsection{Normalized towers can agree to every formal order}
\begin{theorem}[Beyond-all-orders nonuniqueness]\label{thm:ambiguity}
There are infinitely many invariant positive measures $\nu\ne\dd t/t$
for which
\begin{equation}\label{eq:ambiguitynormalizations}
 C_0^\nu(1/2)=\lambda,\qquad K_\nu(0)=K_0.
\end{equation}
For each of them and every fixed $a>0$, the two sequences
$C_m^\nu(a)$ and $C_m(a)$ have the same complete formal expansion in
powers of $2^{-m}$, with the same leading constant.
More quantitatively, they may be constructed so that, for every compact
$I\subset(0,\infty)$, $0<r<4\pi$, and integer $N\ge1$,
\begin{equation}\label{eq:ambiguitybound}
 \frac{|C_m^\nu(a)-C_m(a)|}{K_0\,2^{m(m-1)/2}}
 \le |\eta|M_{I,r}r^{-N}2^{N(N+1)/2-Nm},\qquad a\in I,
\end{equation}
where $0<|\eta|<1$ is a freely chosen perturbation parameter.
The corresponding functions $a\mapsto C_0^\nu(a)$ are nevertheless distinct.
\end{theorem}
\begin{proof}
Choose a nonzero real trigonometric polynomial $p$ of period one satisfying
\begin{equation}\label{eq:twopconditions}
 \int_0^1p(y)\dd y=0,\qquad
 \int_0^\infty\W(t)p(\log_2t)\frac{\dd t}{t}=0,
 \qquad \|p\|_\infty\le1.
\end{equation}
Such choices exist in an infinite-dimensional family. To see this
concretely, take any finite-dimensional span of nonconstant sines and
cosines. Every element already has zero ordinary mean; the weighted
condition imposes at most one linear equation. With two or more
independent modes there is a nonzero solution, which can be rescaled.
Increasing the number of modes gives arbitrarily large families.

Define
\begin{equation}\label{eq:nupperturb}
 \nu(\dd t)=\{1+\eta p(\log_2t)\}\frac{\dd t}{t},
 \qquad0<|\eta|<1.
\end{equation}
This is positive and dyadically invariant. The first zero mean gives
$\nu([1,2))=\lambda$; the second gives $K_\nu(0)=K_0$.
Proposition~\ref{prop:invisible} gives equality of all integer $K$-moments.
Theorem~\ref{thm:universal}, at integer orders, therefore gives the same
formal expansion and leading constant.

For the stronger bound, use the exact renormalization of the difference:
\[
 C_m^\nu(a)-C_m(a)
 =\eta\,2^{m(m-1)/2}\int_0^\infty
       \W(t)p(\log_2t)R_a(2^{-m}t)\frac{\dd t}{t}.
\]
Every term of the Taylor polynomial integrates to zero by
\eqref{eq:allmomentzero}. The global Taylor remainder, $|p|\le1$, and
$K(N)=2^{N(N+1)/2}K_0$ give \eqref{eq:ambiguitybound}.
Finally, the Laplace measures $Q(t)\nu(\dd t)$ are different because
$p$ is nonzero and $Q>0$. Laplace uniqueness implies that the functions
$C_0^\nu$ are different.
\end{proof}

Choosing $N$ as in \eqref{eq:tau}--\eqref{eq:optimal}, with $\rho=m$,
makes the relative difference in \eqref{eq:ambiguitybound} at most
\begin{equation}\label{eq:ambiguityoptimal}
 |\eta|2^{1/8}M_{I,r}\,
       2^{-\frac12(m+\log_2r-1/2)^2}.
\end{equation}
This gives a concrete mathematical meaning to ``invisible to all formal
high-order corrections.'' It does not say that two distinct towers have
the same exact values. Nor does it rule out characterization by a different
asymptotic variable: the next section uses $a\downarrow0$ precisely to
recover the missing information.

\section{A boundary characterization of the canonical tower}\label{sec:uniqueness}

\subsection{A single slope limit selects the invariant measure}
\begin{theorem}[Boundary uniqueness]\label{thm:uniqueness}
Let $f$ be completely monotone on $(0,\infty)$ and satisfy
\eqref{eq:dyadicequation}. Then
\begin{equation}\label{eq:slopecondition}
 \lim_{a\downarrow0}(-a f'(a))=1
\end{equation}
holds if and only if
\begin{equation}\label{eq:canonicalunique}
 f(a)=C_0(a)=\int_0^\infty e^{-at}Q(t)\frac{\dd t}{t}.
\end{equation}
No separate point normalization is needed.
More generally, replacing the limit one by a nonnegative constant $c$
selects $f=cC_0$.
\end{theorem}
\begin{proof}
Let $\nu$ be the invariant representing measure from
Theorem~\ref{thm:classification}. Fix $s>0$ and set $a_N=s2^{-N}$.
Scaling by $2^N$ in the measure gives
\begin{equation}\label{eq:slopescale}
 -a_N f'(a_N)
 =\int_0^\infty st e^{-st}Q(2^Nt)\,\nu(\dd t).
\end{equation}
For each fixed $t>0$, $Q(2^Nt)\to1$.
The dominating kernel $st e^{-st}$ is integrable against $\nu$:
on shells tending to zero it is $O_s(2^j)$, and on shells tending to
infinity it has exponential decay. Each shell has the same finite mass.
Dominated convergence therefore yields
\begin{equation}\label{eq:slopelimit}
 \lim_{N\to\infty}(-a_N f'(a_N))
   =s\int_0^\infty t e^{-st}\,\nu(\dd t).
\end{equation}
If \eqref{eq:slopecondition} holds, the right side equals one for every
$s>0$. Hence the measure $t\nu(\dd t)$ has Laplace transform $1/s$.
This is the Laplace transform of Lebesgue measure $\dd t$ on
$(0,\infty)$, and uniqueness gives $t\nu(\dd t)=\dd t$.
Thus $\nu(\dd t)=\dd t/t$, proving necessity.

For the canonical measure, substitution $u=at$ gives
\[
 -aC_0'(a)=a\int_0^\infty e^{-at}Q(t)\dd t
       =\int_0^\infty e^{-u}Q(u/a)\dd u\longrightarrow1
\]
by $0<Q\le1$ and dominated convergence.
The proof with limit $c$ identifies $t\nu(\dd t)=c\dd t$ in exactly the
same way, including $c=0$.
\end{proof}

Equation \eqref{eq:slopelimit} explains the obstruction in the other
solutions: in general, the dyadic boundary limits retain the phase
information of $\nu$. A full limit as $a\downarrow0$, rather than a
limit along one selected dyadic sequence, removes that information.
The hypothesis genuinely means a limit through \emph{all} positive $a$.

\subsection{A finite-part condition is sufficient}
\begin{corollary}[Finite-part characterization]\label{cor:finitepart}
Within the completely monotone solutions of \eqref{eq:dyadicequation},
the existence of a finite limit
\begin{equation}\label{eq:finitepartcondition}
 \lim_{a\downarrow0}\{f(a)+\log a\}\in\R
\end{equation}
is equivalent to $f=C_0$. For the canonical function, the limit is
\begin{equation}\label{eq:kappa}
 \kappa=-\gamma+
     \int_0^1\frac{Q(t)}t\dd t+
     \int_1^\infty\frac{Q(t)-1}t\dd t,
\end{equation}
where $\gamma$ is Euler's constant.
\end{corollary}
\begin{proof}
Complete monotonicity implies convexity. For $0<h<1$, the two secant
bounds around $a$ give
\begin{equation}\label{eq:convexbounds}
 \frac{f(a)-f((1+h)a)}h
 \le -af'(a)
 \le\frac{f((1-h)a)-f(a)}h.
\end{equation}
Under \eqref{eq:finitepartcondition}, the left and right expressions tend,
as $a\downarrow0$, to $\log(1+h)/h$ and $-\log(1-h)/h$ respectively.
Letting $h\downarrow0$ proves \eqref{eq:slopecondition}, and
Theorem~\ref{thm:uniqueness} applies.

For the converse, split the canonical integral as
\[
 C_0(a)=\int_0^1e^{-at}\frac{Q(t)}t\dd t
       +\int_1^\infty e^{-at}\frac{Q(t)-1}t\dd t
       +\int_1^\infty e^{-at}\frac{\dd t}{t}.
\]
The first two integrals converge to their values at $a=0$ by
Lemma~\ref{lem:Qbounds}. The last is the exponential integral
$E_1(a)=-\log a-\gamma+o(1)$, which gives \eqref{eq:kappa}.
\end{proof}

The finite-part hypothesis is stronger than the slope hypothesis as a
condition on arbitrary functions, but both select exactly the same member
of the dyadic completely monotone family.

\subsection{Uniqueness of the full differential hierarchy}
\begin{corollary}[Canonical gamma-tower characterization]\label{cor:towerunique}
Suppose $f_0,f_1,\ldots$ are completely monotone functions on
$(0,\infty)$ satisfying
\begin{equation}\label{eq:toweraxioms}
 \begin{aligned}
 f_m(a)&=2^m\{f_m(a/2)-f_m((a+1)/2)\} &&(m\ge0),\\
 f_m'(a)&=-f_{m-1}(a) &&(m\ge1),
 \end{aligned}
\end{equation}
and $-af_0'(a)\to1$ as $a\downarrow0$.
Then $f_m=C_m$ for every $m\ge0$. Equivalently, the functions
$\exp((-1)^m m!f_m)$ are exactly the hierarchy in
\eqref{eq:Gammadef}.
\end{corollary}
\begin{proof}
Theorem~\ref{thm:uniqueness} gives $f_0=C_0$.
Inductively, if $f_{m-1}=C_{m-1}$, the differential ladder says that
$f_m-C_m$ is a constant $d_m$.
Subtracting the two level-$m$ dyadic equations gives
$d_m=2^m(d_m-d_m)=0$. This completes the induction.
\end{proof}

This supplies a normalization and a proof of uniqueness in a precise
positive class. It does not establish every possible Weierstrass-product
or complex-analytic assertion one might include in a broader Barnes-type
program. The two issues should be kept separate.

\section{Reproducible computations and what they verify}\label{sec:numerics}

\subsection{Independent integral evaluation}
The archive contains \texttt{verify\_results.py} and two JSON reports.
The coefficient identities were checked exactly with rational symbolic
arithmetic through $k=12$. The numerical calculations were performed
with 48 Gauss--Legendre nodes at 60 decimal digits and repeated with
80 nodes at 80 digits. Both the precision and the quadrature size were
increased; agreement is not based solely on changing the displayed number
of digits.

For $K_0$ and $\Theta$, the program evaluates a single multiplicative
period. Combining \eqref{eq:lognormalW} with $t=2^{y+j}$ gives
\begin{equation}\label{eq:onecellTheta}
 \Theta(\rho)=\lambda\,2^{1/8}
 \int_0^1 A(y)\sum_{j\in\Z}
 2^{-(j+y+\rho-1/2)^2/2}\dd y.
\end{equation}
Periodicity permits reduction of $\rho$ modulo one.
The rapidly convergent Gaussian shell sum makes this inexpensive.

The $C_\rho$ checks are evaluated separately from their \emph{original}
$Q$-integrals, not by inserting the claimed coefficient expansion or the
$R_a$-renormalization. On logarithmic shells the computation uses
\begin{equation}\label{eq:onecellC}
 \begin{split}
 \frac{C_\rho(a)}{2^{\rho(\rho-1)/2}}
  =\lambda\int_0^1\sum_{j\in\Z}
  \exp\bigl(&-\rho\lambda(y+j)-a2^{y+j}\\
       &+\log Q(2^{y+j})-\tfrac\lambda2\rho(\rho-1)\bigr)\dd y.
 \end{split}
\end{equation}
Small factors are evaluated with \texttt{expm1} to avoid cancellation.
The shell values of $Q$ use its own product recurrence. This comparison
can detect an incorrect prefactor or a wrong coefficient in the theorem.

The two runs agree far beyond the digits printed below. For example,
the reported $K_0$ values agree in the first 49 decimal places; their
stored difference is less than $4\cdot10^{-51}$.
The calculations are \emph{not interval-certified}: quadrature stability
and exact symbolic checks support the implementation, while the proofs
establish the theorems. The numerical work does not establish a sharp
constant in the remainder bound or historical novelty.

\subsection{Centered derivative data}
Let
\[
 Z_m=\frac{C_m(1/2)}{K_0\,2^{m(m-1)/2}},\quad
 P_2=1-\frac{2^{-2m}}9,\quad
 P_4=P_2+\frac{248}{2025}2^{-4m},\quad
 P_6=P_4-\frac{4806656}{2679075}2^{-6m}.
\]
Table~\ref{tab:centered} compares successive approximations. For fixed
truncation order the improvement agrees with the even powers predicted by
symmetry.

\begin{table}[htbp]
\centering\small
\begin{tabular}{rcccc}
\toprule
$m$ & $|Z_m-1|$ & $|Z_m-P_2|$ & $|Z_m-P_4|$ & $|Z_m-P_6|$\\
\midrule
4  & $4.322382\,10^{-4}$ & $1.789613\,10^{-6}$ & $7.911730\,10^{-8}$ & $2.782223\,10^{-8}$\\
8  & $1.695392\,10^{-6}$ & $2.850821\,10^{-11}$& $6.353778\,10^{-15}$& $2.031475\,10^{-17}$\\
12 & $6.622738\,10^{-9}$ & $4.350974\,10^{-16}$& $3.799206\,10^{-22}$& $4.918108\,10^{-27}$\\
16 & $2.587007\,10^{-11}$& $6.639065\,10^{-21}$& $2.264533\,10^{-29}$& $1.145315\,10^{-36}$\\
\bottomrule
\end{tabular}
\caption{Absolute errors in the normalized centered expansion. The notation
$x\,10^{-j}$ means $x\times10^{-j}$.}
\label{tab:centered}
\end{table}

At $m=4$, the benefit of further terms is already much less spectacular
than at $m=16$. This is consistent with the fact that the expansion is
divergent: adding indefinitely many terms at a fixed order is not justified.
The full data files also include nonsymmetric shifts and noninteger
orders, for example $a=1$ and $\rho=8.25,8.5,12.75$.

\subsection{The tiny Mellin phase and the visible gamma phase}
The observed relative modulation of $\Theta$ is shown in
Table~\ref{tab:theta}. Its small size is consistent with the exact Fourier
smoothing factor in \eqref{eq:Thetafourier}.

\begin{table}[htbp]
\centering
\begin{tabular}{cc}
\toprule
$\theta$ & $\Theta(\theta)/\Theta(0)-1$\\
\midrule
$1/4$ & $1.289029564918397699\times10^{-18}$\\
$1/2$ & $1.296624062281340296\times10^{-18}$\\
$3/4$ & $7.594497362942597195\times10^{-21}$\\
\bottomrule
\end{tabular}
\caption{Numerical phase variation. Nonconstancy is numerical evidence
here, not an additional theorem assumed elsewhere.}
\label{tab:theta}
\end{table}

For comparison, Table~\ref{tab:phases} gives the explicit extremal
solutions of \eqref{eq:phasefunctions}. Their differences are much easier
to see, although all have $f_\theta(1/2)=\lambda$ and
$f_\theta(1)=\lambda/2$.

\begin{table}[htbp]
\centering
\begin{tabular}{ccc}
\toprule
$\theta$ & $f_\theta(1/3)$ & $-a f_\theta'(a)$ at $a=2^{-24}$\\
\midrule
$0$   & $0.956669839913615728$ & $1.00000628193280958$\\
$1/4$ & $0.956669741884495862$ & $0.99999235648832259$\\
$1/2$ & $0.956665362383835818$ & $0.99999357544426848$\\
$3/4$ & $0.956665460410626991$ & $1.00000750091845969$\\
\bottomrule
\end{tabular}
\caption{Distinct positive normalized solutions. The canonical value is
$C_0(1/3)=0.956667601148143600\ldots$. The slope column is a finite-$a$
evaluation, not a claimed exact limiting value.}
\label{tab:phases}
\end{table}

The program also checks the dyadic equation away from these special
normalization points, using $a=1.3$ and $\theta=0.27$.
The residual is below $10^{-70}$ in the 80-digit run.
Agreement at $a=1$ alone would be an inadequate test because that value
is forced by the normalization.

\subsection{How to reproduce the data}
From the extracted archive, run
\begin{verbatim}
python verify_results.py --dps 60 --nodes 48 --output results_48.json
python verify_results.py --dps 80 --nodes 80 --output results_80.json
\end{verbatim}
The script requires \texttt{mpmath} and \texttt{sympy}, records their
installed versions, and has no network dependency. It checks the
coefficient derivative law, reflection law, and centered parity by exact
symbolic simplification. It does not use numerical differentiation of
$D$ near its zeros, where cancellation would make a direct approach much
less stable.

\section{Interpretation and further research directions}\label{sec:outlook}

The principal mechanism is not cancellation among many alternating
Dirichlet terms. It is a positive product identity. Once $Q$ is paired
with its dyadic uniform-sum counterpart $H$, the resulting multiplicative
cocycle makes its Mellin moments exactly solvable. The difficult
high-order dependence moves into a single analytic function $R_a$ at a
small argument. This is why explicit coefficients and a global remainder
can be obtained together.

The classification theorem gives a complementary lesson. The dyadic
functional equation controls how a measure moves between scales, but it
does not prescribe the distribution of mass within a scale. One-point
normalization determines only that scale's total mass. The all-orders
construction shows that even a full formal expansion at integer orders
can fail to recover the missing phase information. The boundary slope
recovers the whole Laplace transform of the invariant measure, and thus
is strong enough to characterize it.

Several concrete extensions remain worth pursuing. One is to obtain a
closed Fourier or residue formula for the amplitude $A$, and from it a
fully explicit proof and bound for the nonconstant modulation of $\Theta$.
Another is to determine the leading exponentially small discrepancy in
Theorem~\ref{thm:ambiguity}, rather than only its Gaussian-in-order upper
bound. The poles of $R_a$ at imaginary multiples of $4\pi$ and the chosen
phase measure should enter such a refinement, but this article does not
supply that additional expansion.

A separate quantitative task is to replace the tube supremum $M_{I,r}$ by
simple computable constants optimized for a given shift interval.
The proof of Lemma~\ref{lem:tubebound} already separates a compact region
from a controlled tail, so interval arithmetic could turn it into a
numerical certificate. The present tables should not be mistaken for that
certificate.

The positive-measure classification and the proof of the uniform remainder
also have clear formalization targets: product convergence, pushforward
identities, the Mellin recurrence, Cauchy estimates on a tube, and Laplace
uniqueness. None of these targets requires asserting that the results
already appear in the repository's Lean code. Finally, analogous product
pairs for other dilation factors may produce related invariant-measure
classifications; a radix change must be made at the level of the product
equation rather than by replacing every two in the final formulas without
checking the underlying sequence.

\appendix
\section{A compact audit of hypotheses and conclusions}\label{app:audit}

The asymptotic parameter is real $\rho\to\infty$, with $a$ in a fixed
compact subset of $(0,\infty)$. Uniformity up to $a=0$ is not claimed:
indeed, $R_0(x)=1/H(x)$ is unbounded as $x\to\infty$, which invalidates
the tube supremum used in the proof. The restriction is structural rather
than cosmetic.

The all-orders series has zero convergence radius as a series in
$h=2^{-\rho}$. A finite expansion with a proved remainder is the intended
meaning of every $\sim$ symbol in this article. Optimal truncation gives
an upper bound of Gaussian size in $\rho$, but not a matching lower bound.
The finite coefficient formula is exact independently of numerical tests.

The uniqueness theorem assumes complete monotonicity, not merely
positivity of $f$ or of an exponentiated gamma function. This assumption
makes a positive Laplace measure available and is essential to the stated
classification proof. The normalized nonunique examples satisfy this
stronger condition, so they are not counterexamples obtained by leaving
the proposed positive class.

The ambiguity theorem compares integer-order expansions. In general its
perturbations need not preserve the entire function
$\Theta_\nu(\rho)$ at noninteger orders. Equality of every integer Mellin
moment is weaker than equality of the full Mellin transform, just as the
explicit Stieltjes family demonstrates.

The term ``gamma tower'' refers exactly to
$\exp((-1)^m m!C_m(a))$, or its measure-deformed counterpart. The article
proves a real positive-class characterization; it does not silently assume
a preferred continuation to all complex $a$, a canonical Weierstrass
factorization, or an inversion transseries for that tower.

\section{Source and novelty audit}\label{app:sources}

The supplied repository document was read as the public raw LaTeX source
at the URL in reference~\cite{Atlas}; its consolidated document is dated
August 2026. The comparison used its Mellin-continuation material, its
automatic gamma-hierarchy material, and its stated frontier on hierarchy
normalization. A moving repository URL is not an immutable version pin;
no unverified commit hash is assigned to it here.

Classical Dirichlet-series facts were checked against the sources listed
below. The primary text of the complete-monotonicity representation was
available in~\cite{AJ}; the present measure equation and its classification
are derived in Section~\ref{sec:classification}, rather than attributed to
that reference. The general historical statements about Mellin
log-periodicity and lognormal moment indeterminacy are background only.
None of those references is being credited with the particular uniform
remainder theorem or the boundary argument without evidence.

The 2001 Alkauskas preprint was located on its author's bibliography, but
the linked PDF did not load during this investigation. It is recorded
because of its potentially relevant subject, not cited as a text that was
fully examined. Consequently, historical priority remains a separate
verification task. The claims of this article are the stated mathematical
theorems and their proofs, supplemented by reproducible but non-interval
numerical checks.

\clearpage
\begin{thebibliography}{99}
\bibitem{Atlas}
V.~Reshetnikov, \emph{The Thue--Morse Sequence: Formula Atlas and
Fabius--Rvachev Frontier Results}, ProveIt repository, consolidated
August 2026, public source consulted 4 September 2026.\\
\url{https://github.com/VladimirReshetnikov/ProveIt/tree/main/Analysis/FabiusFunction/docs/semi-formalized-research-frontiers/drafts/thue-morse/Thue_Morse_Atlas_and_Frontiers}

\bibitem{AC}
J.-P.~Allouche and H.~Cohen, Dirichlet series and curious infinite products,
\emph{Bulletin of the London Mathematical Society} \textbf{17} (1985),
531--538. \url{https://doi.org/10.1112/blms/17.6.531}.

\bibitem{AMP}
J.-P.~Allouche, M.~Mend\`es France, and J.~Peyri\`ere,
Automatic Dirichlet series, \emph{Journal of Number Theory}
\textbf{81} (2000), 359--373.
\url{https://doi.org/10.1006/jnth.1999.2487}.

\bibitem{Allouche2015}
J.-P.~Allouche, Thue, combinatorics on words, and conjectures inspired by
the Thue--Morse sequence, \emph{Journal de Th\'eorie des Nombres de Bordeaux}
\textbf{27} (2015), 375--388.
\url{https://arxiv.org/abs/1401.3727}.

\bibitem{Arias}
J.~Arias de Reyna, Arithmetic of the Fabius function,
\emph{Integers} \textbf{18} (2018), Paper A51.
\url{https://math.colgate.edu/~integers/s51/s51.pdf}.

\bibitem{FGD}
P.~Flajolet, X.~Gourdon, and P.~Dumas,
Mellin transforms and asymptotics: harmonic sums,
\emph{Theoretical Computer Science} \textbf{144} (1995), 3--58.
\url{https://specfun.inria.fr/dumas/Publications/FlGoDu95.pdf}.

\bibitem{Heyde}
C.~C.~Heyde, On a property of the lognormal distribution,
\emph{Journal of the Royal Statistical Society, Series B}
\textbf{25} (1963), 392--393.
\url{https://doi.org/10.1111/j.2517-6161.1963.tb00521.x}.

\bibitem{AJ}
R.~Aguech and W.~Jedidi, New characterizations of completely monotone
functions and Bernstein functions, a converse to Hausdorff's moment
characterization theorem, \emph{Arab Journal of Mathematical Sciences}
\textbf{25} (2019), 57--82. Theorem~1 states the positive Laplace
representation and uniqueness used here.\\
\url{https://ajms.ksu.edu.sa/sites/ajms.ksu.edu.sa/files/imce_images/new-characterizations-of-completely-monotone-functions-a_2019_arab-journal-o.pdf}

\bibitem{Toth}
L.~T\'oth, \emph{Linear Combinations of Dirichlet Series Associated with the
Thue--Morse Sequence}, arXiv:2211.13570 (2022).
\url{https://arxiv.org/abs/2211.13570}.

\bibitem{Alkauskas}
G.~Alkauskas, \emph{Dirichlet series associated with Thue--Morse sequence},
Vilnius University, preprint 2001-15 (2001). Bibliographic record located
on the author's page; linked full text not retrieved in this investigation.\\
\url{https://web.vu.lt/mif/g.alkauskas/math/index.html}.
\end{thebibliography}
\end{document}
